% Pista ana-24j-q1 · Anàlisi
% Procedència: PAU Catalunya 2024 · Convocatòria de juny · Sèrie 1 · Qüestió 1
\documentclass[11pt,a4paper]{article}
\usepackage{pista}
\pistainfo{Catalunya 2024}{Convocatòria de juny}{Sèrie 1}

\begin{document}

\section*{\textcolor{blauPista}{Qüestió 1}}

\noindent Considereu la funció
\[
f(x) = 2\,\dfrac{\ln x}{x},
\]
definida per a $x > 0$.

\subsection*{\textit{a)} \normalfont Estudieu-ne els màxims i els mínims, i les zones de creixement i de decreixement. \hspace*{\fill}\textnormal{[1~punt]}}

\pista{En primer lloc, calcula $f'(x)$ amb la regla de la derivada d'un quocient:
\[
\left(\frac{u}{v}\right)' = \frac{u' \cdot v - u \cdot v'}{v^{2}}
\]
A continuació, resol $f'(x) = 0$. En aquest cas, només cal igualar a zero el numerador. Recorda que $\ln e = 1$.

\medskip
Finalment, fes una taula de signes de la funció derivada $f'(x)$, però només per a $x > 0$, perquè t'has de restringir al domini de la funció $f(x)$. El denominador sempre és positiu, perquè és un quadrat. Això vol dir que el signe de la funció $f'$ depèn només del numerador. En el punt on una funció passa de créixer a decréixer hi ha un màxim relatiu; en canvi, on passa de decréixer a créixer, hi ha un mínim relatiu.}

\subsection*{\textit{b)} \normalfont Aquesta funció té asímptotes? Feu un esbós de la seva gràfica. \hspace*{\fill}\textnormal{[1~punt]}}

\pista{Tens gairebé tota la feina feta a l'apartat \textit{a)}. Et falten els límits als dos extrems del domini:
\begin{itemize}
\item \textit{Quan $x \to 0^{+}$.} El numerador, $2 \ln x$, tendeix a $-\infty$, i el denominador, $x$, tendeix a $0$ amb valors positius. No és una indeterminació: quin signe té el quocient? Si el límit és $+\infty$ o $-\infty$, hi ha una asímptota vertical a $x = 0$.
\item \textit{Quan $x \to +\infty$.} El numerador i el denominador tendeixen a $+\infty$, però $x$ creix molt més de pressa que $\ln x$. Prova valors grans: $f(10) \approx 0{,}46$, $f(100) \approx 0{,}092$ i $f(1000) \approx 0{,}014$. Cap a quin valor s'acosten els valors de la funció $f(x)$? Si el límit és un nombre real (finit), llavors significa que hi ha una asímptota horitzontal.
\end{itemize}
Per a l'esbós, fes servir aquests límits, el creixement i el decreixement de l'apartat \textit{a)}, i el punt on la gràfica talla l'eix d'abscisses, que trobaràs resolent $f(x) = 0$. No et cal la segona derivada.}

\subsection*{\textit{c)} \normalfont Calculeu l'equació de la recta tangent a la gràfica de $y = f(x)$ en el punt d'abscissa $x = 1$. \hspace*{\fill}\textnormal{[0,5~punts]}}

\pista{El pendent de la recta tangent que busques és $f'(1)$, i el punt és $\bigl(1, f(1)\bigr)$. Per tant:
\[
y - f(1) = f'(1)\,(x - 1)
\]
Ja tens $f'(x)$ de l'apartat \textit{a)}: només cal que calculis $f(1)$ i $f'(1)$, i que els substitueixis a la fórmula. Recorda que $\ln 1 = 0$.}

\end{document}
